\section{Conclusão}
A pergunta-chave foi respondida por cinco decisões verificáveis no código e no log: a mesma função gera a tela humana e a observação do agente (paridade de signos); a MoLIC restringe a resposta do LLM às falas $u{:}$ do estado; o comportamento da persona sai de mecanismos sorteados com semente, fora do prompt, e cada intervenção é registrada; a ajuda chega só por fala registrada e estratificada por tipo de dúvida; e o sucesso é decidido por um verificador que lê apenas o estado final. Essas escolhas tornaram auditável cada diferença entre o que a tela mostrou e o que o agente recebeu, e permitiram ao piloto revelar defeitos de protocolo (chave ambígua em C2, R3 inoperante, confusão entre exploração e ajuda) antes de qualquer uso das tarefas de validação ou de estudo com pessoas. Os próximos passos são revisar C2 e R3 na calibração, congelar $\theta$, prompt e modelo, e só então executar V1 e V2. O agente continua um instrumento de depuração e não substitui participantes \cite{lu2025uxagent}.
